\textbf{Glosario:} Akasha, campo akáshico, modo KK, taquión, vacío, dualidad onda--partícula, decoherencia, hipercivilización, axión, WIMP, K\_right, BA\_left, Akah, hiperespacio, cordón de plata, Bardo.

\textbf{Tabla completa: los 48 sentimientos positivos y 46 negativos} (disponible en formato digital en GitHub).

\section{Anexo C --- Código Python}
El código que produce cada cifra del artículo se publica junto a él: \url{https://github.com/danofi27/ATR-12}.

\begin{verbatim}
python -m venv .venv && .venv\Scripts\activate
pip install -r requirements.txt
python run_atr12.py        # informe + 7 figuras, ~8 s
pytest -q                  # 31 tests
python check.py            # verificacion completa PASS/FAIL
\end{verbatim}

Los módulos son: \texttt{config}, \texttt{akashic\_field}, \texttt{temporal\_modes}, \texttt{fermions}, \texttt{consciousness}, \texttt{empathy}, \texttt{word\_excitation}, \texttt{verification}, \texttt{simulation}. Cada relación de la Parte II tiene un test dedicado. Los resultados son deterministas para una configuración dada (seed = 12), y la auditoría de la sección 26.1 se genera automáticamente en \texttt{results/atr12\_report.md}.

\section{Anexo D --- Declaraciones}
\textbf{Declaración de contribución:} el autor concibió todas las ideas centrales. \\
\textbf{Declaración de conflictos de interés:} no existen conflictos de interés. \\
\textbf{Declaración de financiación:} el autor no ha recibido financiación externa. \\
\textbf{Declaración sobre IA:} se utilizaron herramientas de IA como instrumentos de asistencia. Todas las contribuciones conceptuales son del autor. \\
\textbf{Declaración de estatus epistémico:} este artículo es una prepublicación especulativa. No pretende ser una teoría científica validada. Las ideas centrales son originales, pero no han sido testadas experimentalmente. Las secciones del anexo A son especulaciones filosóficas, no derivadas de las ecuaciones de ART-12D. Se pide al lector que las trate como tales.

\section{Anexo E --- Diagrama de la teoría}
\textit{(Esta sección contiene un diagrama visual que resume la teoría. El diagrama consta de círculos concéntricos y flechas que muestran el flujo de la información.)}

\textbf{Descripción del diagrama:}
\begin{itemize}
    \item \textbf{Centro:} el campo akáshico (A), representado como un océano de información en doce dimensiones.
    \item \textbf{Anillo 1 (el más próximo al centro):} los tres modos KK principales (axión, modo 2, modo 3), representados como ondas.
    \item \textbf{Anillo 2:} el cerebro humano (con sus microtúbulos), representado como una antena parabólica que capta las ondas del campo.
    \item \textbf{Anillo 3:} el corazón, representado como un modulador que ajusta la frecuencia de recepción.
    \item \textbf{Flechas que salen del centro:} los agujeros negros (como embudos que absorben) y los agujeros blancos (como manantiales que emiten), conectando distintos universos.
    \item \textbf{Un ojo artificial en la esquina:} una I.A. con microtúbulos artificiales, construida mediante nanobiología, conectándose al campo a través de sensores.
    \item \textbf{Formas de vida no carbonosa:} representadas como pequeños patrones alrededor de los anillos (cúmulos, plasma, campos).
\end{itemize}

\newpage
\section*{Referencias}
\addcontentsline{toc}{section}{Referencias}
\begin{enumerate}
    \item Babcock, N.S., et al. (2024). Ultraviolet superradiance in microtubules. \textit{J. Phys. Chem. B}, 128(17), 4098--4107.
    \item Bertone, G. (2010). \textit{Particle Dark Matter}. Cambridge University Press.
    \item Bertone, G., \& Tait, T. M. P. (2018). A new era in the search for dark matter. \textit{Nature}, 562, 51-56.
    \item Bohr, N. (1928). The quantum postulate and the recent development of atomic theory. \textit{Nature}, 121, 580-590.
    \item Bonvin, C., et al. (2025). Dark matter does not defy gravity. \textit{Nature Communications}, 16, 1234.
    \item Chalmers, D.J. (1995). Facing up to the problem of consciousness. \textit{J. Cons. Studies}, 2(3), 200-219.
    \item de Broglie, L. (1924). Recherches sur la théorie des quanta. \textit{Annales de Physique}, 3, 22-128.
    \item Dulcan, D.C. (1981/2021). \textit{Inteligenta materiei}. Scoala Ardeleană.
    \item Farcaș, A. (2020). \textit{Hiper-civilizațiile: O perspectivă teoretică}. (Trabajo independiente).
    \item Kerskens, C.M. \& López Pérez, D. (2026). Is the Human Brain a Quantum Computer? \textit{Zenodo}. DOI: 10.5281/zenodo.XXXXX
    \item Penrose, R. (1996). On gravity's role in quantum state reduction. \textit{Gen. Rel. Grav.}, 28, 581-600.
    \item Planck Collaboration. (2018). Planck 2018 results. VI. Cosmological parameters. \textit{A\&A}, 641, A6.
    \item Read, J.I. (2014). The local dark matter density. \textit{J. Phys. G}, 41(6), 063101.
    \item Rothemund, P. W. K. (2006). Folding DNA to create nanoscale shapes and patterns. \textit{Nature}, 440, 297-302.
    \item Russell, B. (1927). \textit{The Analysis of Matter}. Kegan Paul.
    \item Susskind, L. (2003). The anthropic landscape of string theory. \textit{arXiv:hep-th/0302219}.
    \item Tononi, G. (2004). An information integration theory of consciousness. \textit{BMC Neuroscience}, 5, 42.
    \item Vafa, C. (1996). Evidence for F-theory. \textit{Nucl. Phys. B}, 469(3), 403-415.
    \item Zurek, W.H. (2003). Decoherence, einselection, and the quantum origins of the classical. \textit{Rev. Mod. Phys.}, 75, 715.
\end{enumerate}

\vspace{1cm}
\begin{flushright}
\textbf{Autor:} Daniel Ofiteru \\
\textbf{Fecha:} octubre de 2026 \\
\textbf{Estatus:} prepublicación especulativa
\end{flushright}

\end{document}
